\section{この結果で言えないこと}
\label{sec:limits}

\subsection{解像度の限界（第 \ref{sec:validation} 章より）}

第 \ref{sec:validation} 章の交差検証から、読み取ってはいけない範囲がはっきりしている。

\begin{itemize}
  \item \textbf{日単位の傾向は無い。} 3 領域とも、1 日ビンは最良より
        6 標準誤差以上悪い。本統計に日単位の集計を置いていないのはそのためである。
  \item \textbf{北海道の半月ごとの違いは読み取らない。}
        月単位と統計的に区別できない（$+0.5$ 標準誤差）。
  \item \textbf{冬の半月ごとの違いは読み取らない。}
        冬の中では、時期を細かく知っても落下位置の予測が改善しない（誤差削減 0 \%）。
  \item \textbf{地上付近にこの統計を使わない。}
        \SI{850}{\hecto\pascal} の誤差削減は 4--6 \% しかない。
        打ち上げ当日の地上風は予報で判断すること。
  \item \textbf{気候値は予報の代わりにならない。}
        最良の粒度でも説明できるのは落下位置の分散の 70 \%（和歌山）、
        72 \%（愛媛）、30 \%（北海道）で、残りは当日の天気である。
        風そのものでみても、季節が説明するのは分散の 45--47 \%（和歌山・愛媛）、
        18 \%（北海道）にとどまる
        （第 \ref{sec:fidelity} 章）。
\end{itemize}

なお「半月単位の傾向がある」という言い方自体が正確ではない。
実在するのは 1--2 か月スケールの滑らかな季節変化で、
半月ビンはそれを階段で近似したものである（第 \ref{sec:validation} 章）。

\subsection{予報の誤差は測っていない（第 \ref{sec:fidelity} 章より）}

第 \ref{sec:fidelity} 章はシミュレータが風を固定していることの誤差を測ったが、
測れるのは 4 つある誤差源のうち 3 つだけである。
\textbf{与えた予報そのものが外れている分は測れていない。}
それには「過去に出された予報」の保管庫が要るのに対し、
本研究が持っているのは「過去に実際どうだったか」の再解析だからである。

したがって第 \ref{sec:fidelity} 章の \SIrange{2.8}{5.3}{\kilo\metre} という数字は、
\textbf{予報が完璧だった場合の下限}である。
実際の運用では予報誤差が上乗せされ、これより外れる。
実務ではそちらのほうが大きい可能性が高い。

\subsection{落下位置には地図因子の系統誤差が含まれる（第 \ref{sec:fidelity} 章より）}

第 \ref{sec:hokkaido} 章から第 \ref{sec:ehime} 章の落下位置は
$\Delta x = \sum_k u_k T_k$ という閉形式で計算している。
この式には緯度が現れないので、飛行中に機体が南北へ動いて
経度 1 度あたりの距離が変わる効果を持っていない。

第 \ref{sec:fidelity} 章で球面の幾何を正しく扱う積分器と突き合わせたところ、
和歌山の冬でずれは中央値 \SI{0.44}{\kilo\metre}、
10 年ぶんの最大で \SI{3.45}{\kilo\metre} であった。
ビンどうしの比較には影響しない大きさだが、
1 回の飛行の予測に使う場合はこの分を見込む必要がある。

\subsection{愛媛の領域はデータ範囲の西端に接している}

取得したデータの範囲（\ang{132}--\ang{150}E）は、愛南町の西 \SI{51}{\kilo\metre} で終わる。
そのため愛媛の領域は打ち上げ地点を中心にできず、東へ伸ばしてある
（第 \ref{sec:ehime} 章）。

影響は 2 つある。
第 \ref{sec:fidelity} 章の軌道計算では、愛媛から上げた軌道の \SI{5.9}{\percent} が
西端からデータの外へ出て、集計から外れた。
外れたのは夏に西へ流れた短い飛行が中心で、その偏りは誤差の中央値で 2 \% 未満である。
また、豊後水道より西（九州側）の風は本報告のどの図にも入っていない。
夏の西向きの飛行について、九州寄りの風の違いを確かめるには範囲を西へ広げた取り直しが要る。

\subsection{再解析は観測ではない}

第 \ref{sec:data} 章で述べたとおり、JRA-3Q の値は観測に強く拘束されたモデル推定である。
日本の陸上の対流圏では観測が密なので実態に近いが、
外洋上や高度 \SI{30}{\kilo\metre} 以上では、
その場を直接測ったものは少なく、モデルの寄与が相対的に大きい。
本報告で最もスペースバルーンに関係する高度帯（\SIrange{25}{35}{\kilo\metre}）は、
ちょうどその「観測が薄い」側にある。

\subsection{解像度が地形を解いていない}

\ang{0.375} は約 \SI{40}{\kilo\metre} である。
日高山脈も紀伊山地も四国山地も、この格子では 2--3 点でしか表現されない。
したがって本報告は\textbf{谷筋の風や斜面下降風のような地形性の現象を一切含まない}。
打ち上げ地点そのものの地上風は、この統計ではなく現地の観測で判断すべきである。

上空数キロ以上では地形の影響が薄れるので、
スペースバルーンの飛行の大部分についてはこの制約は緩い。

\subsection{10 年は経年変動をならすには短い}

成層圏の東西風には準 2 年周期振動（QBO）があり、
エルニーニョ・南方振動は対流圏の循環を年ごとに動かす。
10 年（QBO でいえば約 4--5 周期）では、これらの位相が完全にはならされない。
あるビンの値が「その時期のふつう」なのか
「たまたまこの 10 年でそうだった」のかは、本報告では分離できない。

より長い期間を使えばこの点は改善するが、
今度は観測網の変化（衛星観測の導入など）が入ってくる。
どちらを取るかの判断であって、片方が正しいわけではない。

\subsection{気圧面上の統計であって等高度面ではない}

統計は気圧面ごとに取っている。
ある気圧面の実際の高度は季節や天気で数百メートル動くので、
「高度 \SI{20}{\kilo\metre} の統計」と読み替えるときには、
その高度が平均値であることを意識する必要がある。
高度の標準偏差も出力に含めてあるので、必要ならそちらを見てほしい。

\subsection{半月の平均は「その日」を教えてくれない}

これは気候値というものの性質そのものである。
「3 月前半の平均風は西 \SI{30}{\metre\per\second}」と分かっても、
来年の 3 月 10 日がそうだとは限らない。
分位点と分散楕円を併記しているのは、
平均からどれだけずれうるかまで含めて計画に使えるようにするためである。
実際に飛ばす日の判断には、この統計ではなく予報を使うこと。
